\section*{Read this first}
\addcontentsline{toc}{section}{Read this first}

This part is written to be read in one sitting, after the one-month pause,
before diving back into code. It answers four questions: what is this
project, why does it look the way it does, what has actually been
accomplished, and what should happen next.

\subsection*{What RoboReason is}

RoboReason ROS2 lets an operator type (or eventually speak) a natural-language
instruction -- ``pick the red cube and place it on the tray'' -- and have a
UR5cb arm with an OnRobot RG2 gripper carry it out. A language model (LLM) or
vision-language model (VLM) turns the instruction plus a description of the
scene into a structured sequence of primitive skills (\texttt{approach},
\texttt{pick}, \texttt{release}, \texttt{move\_home}, \texttt{wait}), which a
plan manager validates and a skill executor runs on the real (or simulated)
arm. A FastAPI web GUI supervises the whole ROS2 stack, the camera, and the
UR driver, and gives an operator a chat-style interface with live execution
feedback, an emergency stop, and per-run debug capture.

\subsection*{Why LLM/VLM and not VLA}

\textbf{The project did not start here.} The \texttt{docs/} folder in the
adjacent working directory
(\texttt{Robot\_AI\_...Technical\_Steps}, \texttt{Vision\_Language\_Action\_
(VLA)\_Models\_in\_Robotics}, \texttt{VLA\_Future\_Directions}) documents an
earlier direction built around end-to-end Vision-Language-Action models --
networks trained to map pixels and language directly to robot actions. That
approach was \textbf{deliberately abandoned} in favor of the current
LLM/VLM-plus-classical-planning architecture. Those older documents are kept
only as historical record; nothing in them describes the system as it exists
today, and they should not be used as a reference for current work.

The practical reasons for the pivot, as reflected in the codebase's design:
\begin{itemize}[nosep]
  \item \textbf{Interpretability and debuggability.} An LLM/VLM plan is a
    readable JSON list of named skills with explicit coordinates -- every
    step can be logged, inspected, and corrected. A VLA's action output is a
    black box by comparison, which makes hardware debugging (the dominant
    cost on this project, see Part~\ref{part:technical}) far harder.
  \item \textbf{No training data / no fine-tuning loop required.} Off-the-shelf
    foundation models (Groq, Nebius, OpenAI, Anthropic, Gemini) are used
    directly through their APIs, with prompt engineering as the main lever --
    no dataset collection or model training infrastructure needed.
  \item \textbf{Deterministic geometry where it matters.} As the benchmark
    study in Section~\ref{sec:benchmark} shows, most real-world failures on
    this project trace back to \emph{arithmetic and reasoning} mistakes by
    the model (misjudging heights, offsets, or which object a value maps to),
    not to perception. The current architecture lets those computations be
    pulled out of the model and done deterministically in Python
    (Section~\ref{sec:tech-decisions}), which is not possible with an
    end-to-end VLA.
\end{itemize}

\subsection*{Timeline at a glance}

The project's git history runs from the initial commit on
\textbf{27 May 2026} to the most recent committed work on
\textbf{22 July 2026} -- roughly eight weeks, organized into five working
sessions of increasing scope. Figure~\ref{fig:exec-timeline} gives the
high-level shape; Section~\ref{sec:timeline} gives the full session-by-session
account.

\begin{figure}[H]
\centering
\begin{tikzpicture}[
  every node/.style={font=\small},
  milestone/.style={rectangle, rounded corners, draw=rrblue, fill=rrblue!8,
    minimum width=2.6cm, minimum height=1.1cm, align=center, text width=2.5cm},
  dateline/.style={rrgray, font=\scriptsize}
]
  \draw[-{Latex[length=2mm]}, thick, rrblue!60] (-0.3,0) -- (14.8,0);

  \foreach \x/\lab/\date in {
    0/{VLA pivot \\ \& bring-up}/27 May,
    3/{VLM pipeline \\ \& calibration}/mid-Jun,
    6/{Web GUI \\ (6 phases)}/mid-Jun,
    9/{VLM$\to$LLM \\ hybrid}/late Jun,
    12/{Grasp geometry \\ hardening}/early Jul,
    14.5/{Benchmark study \\ (120 trials)}/15--20 Jul
  } {
    \draw[rrblue, thick] (\x,0.08) -- (\x,-0.08);
    \node[dateline, below=2pt] at (\x,-0.1) {\date};
  }

  \node[milestone] at (0,1.3) {Session 1\\Bring-up};
  \node[milestone] at (3.5,1.9) {Session 2\\Code review \& refactor};
  \node[milestone] at (7.5,1.3) {Session 3\\Web GUI + e-stop};
  \node[milestone] at (11,1.9) {Session 4\\VLM$\to$LLM hybrid};
  \node[milestone] at (14.5,1.3) {Session 5\\Benchmark};

  \draw[rrblue!40] (0,1.05) -- (0,0.15);
  \draw[rrblue!40] (3.5,1.55) -- (4.5,0.15);
  \draw[rrblue!40] (7.5,1.05) -- (7,0.15);
  \draw[rrblue!40] (11,1.55) -- (10.5,0.15);
  \draw[rrblue!40] (14.5,1.05) -- (14.5,0.15);
\end{tikzpicture}
\caption{Project timeline, 27 May -- 22 July 2026 (five working sessions).}
\label{fig:exec-timeline}
\end{figure}

\subsection*{What has actually been delivered}

\begin{itemize}[nosep]
  \item A working \textbf{LLM planning pipeline} (\texttt{mode=LLM}) driven
    by a static hand-authored scene description, and a working
    \textbf{VLM planning pipeline} (\texttt{mode=VLM}) driven by a real
    camera with ChArUco-calibrated pixel-to-3D deprojection, plus a
    \textbf{hybrid VLM$\to$LLM mode} (\texttt{mode=VLM\_LLM}) that grounds a
    scene with the VLM and then plans over it with the LLM.
  \item Six interchangeable \textbf{reasoning methods}
    (FHP/FFHP, ReAct, Chain-of-Thought Self-Consistency, Self-Refine,
    Tree-of-Thought, Always-Act), all supporting both point and bounding-box
    pixel grounding for VLM modes.
  \item A full \textbf{web GUI} (\texttt{robo\_reason\_gui}) that supervises
    the ROS2 stack, the UR driver (with auto-retry against a known flaky
    startup), and the camera service as child processes, with live chat,
    execution streaming, an emergency stop, and per-run debug capture.
  \item A series of \textbf{hardware-driven correctness fixes} to the
    grasp/release geometry (Section~\ref{sec:tech-decisions}), each traced to
    a specific observed failure on the real robot and verified against
    captured debug data, not just unit tests.
  \item A \textbf{120-trial benchmark} directly comparing the LLM and VLM
    pipelines across six task conditions, with a reusable measurement
    harness (Section~\ref{sec:benchmark}) that produces publication-ready
    figures.
\end{itemize}

\subsection*{Headline result}

Across the benchmark's 120 trials, the \textbf{LLM pipeline} was safer and
more successful than the \textbf{VLM pipeline} (95\% task-safety / 97\%
task-success vs.\ 82\% / 66\%). The more interesting finding, reached by
reading the operator's free-text failure notes rather than trusting the
aggregate numbers alone, is that VLM's dominant failure mode is \emph{not}
color perception -- it is \emph{arithmetic/value-mapping reasoning} under the
harder tasks. Section~\ref{sec:benchmark} gives the full breakdown.

\subsection*{State of the working tree}

As of this report, all Session 5 work (deterministic release-geometry fixes,
grasp-width sampling fix, the full \texttt{benchmark/} harness, and
\texttt{docs/TODO.md}) has been committed by the operator, so the repository
reflects the state described in this report with no dangling uncommitted
changes carried across the pause.

\subsection*{Picking this back up in September}

Recommended first session, in order:
\begin{enumerate}[nosep]
  \item \textbf{Re-verify hardware, don't assume.} The two Session 5
    geometry fixes (\texttt{\_fix\_object\_height},
    \texttt{\_fix\_release\_height}) were validated against \emph{replayed}
    debug data, not a fresh end-to-end run. Re-run the \texttt{sort\_hard}
    and \texttt{arith\_hard} prompts that originally surfaced the bugs they
    fix, and confirm no regressions, before trusting them for new work.
  \item \textbf{Check the UR driver and camera calibration still hold.} Both
    are physical-setup-dependent (Section~\ref{sec:open-issues}); a month of
    disuse is enough for a cable, mount, or lighting change to drift the
    ChArUco calibration.
  \item \textbf{Read Section~\ref{sec:roadmap}} for the prioritized backlog --
    it reflects everything in \texttt{docs/TODO.md} plus items raised but
    deliberately deferred during the benchmark push (notably the
    bounding-box scene schema change, item R4).
\end{enumerate}

The rest of this report (Part~\ref{part:technical}) is the full technical
record: architecture, the complete session-by-session history, every
significant bug fix with its root cause, the benchmark methodology and
results in detail, the current open-issues list, and the prioritized
roadmap.
